\section{One rule as a quantum computer}\label{sec:intro}
A fixed local rule can contain a universal processor while leaving the desired computation to its initial state. Under continuing noise this requires more than protecting quantum data: addresses, program symbols, clocks, routing decisions and every bit of classical correction work must also survive. We construct such a rule in one spatial dimension. Its local register, gate range and elementary period are fixed once. Larger computations and smaller errors change only the length of the ring, the product preparation and the stopping time.

Here and below the error of a probability law is total variation distance. The quantum endpoint result uses one half of the trace norm, including a reference. A period means one complete application of the fixed elementary circuit. The elementary gates include reset to zero, and noise follows every layer, also on idle and reset wires. Section~\ref{sec:model} specifies the channels exactly.

\begin{maintheorem}[Product preparation and classical readout]\label{thm:main}
There exist a single finite-range, translation-invariant, time-periodic quantum circuit on rings, a finite number $Q_{\rm site}$ of qubits per site, and constants $p_0>0$ and $A\ge1$ with the following property. For every $n\ge1$, $D\ge0$, every depth-$D$ Toffoli/Hadamard circuit $C$ on $n$ qubits and every $0<\varepsilon<1$, use a sufficient integer precision request $h$ as specified in Section~\ref{sec:precision}. There is a polynomial-time computable product preparation in which only classical basis registers carry nonzero data and every quantum register is zero. The classical state describes $C$ and $h$.

For every $0<p\le p_0$, apply independent single-qubit depolarization of rate $p$ to every wire after every elementary layer. After $t$ periods on $L$ sites, measure all qubits in the computational basis. A deterministic polynomial-time classical decoder depending only on $n,D,h$ gives the measured output law of $C\ket{0}^{\otimes n}$ within total variation distance $2^{-h}\le\varepsilon$. In particular the decoder does not read or evaluate the gate list of $C$. Both resources satisfy
\begin{equation}\label{eq:resource}
 L,t\le A\bigl(n+D+\log_2(1/\varepsilon)\bigr)^A.
\end{equation}
The hardware, preparation, decoder and stopping time do not depend on $p$. Explicit dependence of $p_0$ and $A$ on the fixed construction is given in \eqref{eq:threshold} and \eqref{eq:A}.
\end{maintheorem}

D\"unnweber, Styliaris and Trivedi~\cite[Theorem~1]{DST2026} use a fixed local update channel in an alternating two-dimensional brickwork schedule, with noisy control and an initial-state program. For a sufficiently correct hierarchical input they prove a positive local fault-path threshold, using $O(N\operatorname{polylog}(ND/\delta))$ qudits and $O(D\operatorname{polylog}(ND/\delta))$ steps after translation-invariant logical compilation. Their Section~VI describes encoded preparation in polylogarithmic depth. Theorem~\ref{thm:main} instead supplies a computable basis-product start on a line, with polynomial space and time and the independent noise models of Section~\ref{sec:precision}. Theorem~\ref{thm:general} below extends our rule to their local fault-path noise convention, including coherent errors and bath correlations obeying the specified superset bound. Their polylogarithmic overhead remains stronger than our polynomial bound. They explicitly suggested embedding a quantum code in G\'acs's one-dimensional hierarchy as a route to the one-dimensional case. This suggestion already appears in their January 2026 manuscript and is retained in its later versions. We follow that route.

One-dimensional fault-tolerant quantum computation with local gates and no intermediate measurements or classical processing already appears in Aharonov and Ben-Or's threshold theorem~\cite[Theorem~13]{ABO2008}. Gottesman~\cite{Gottesman2000} gives low-dimensional local-gate thresholds, with next-to-nearest-neighbor gates in one dimension, and suggests interspersing reliable classical regions of the G\'acs type. Those constructions have polylogarithmic overhead, but they compile a target circuit into a spacetime gate list. Here the spacetime rule is fixed and periodic, and only its initial state contains the target program. The interpretation of the Aharonov--Ben-Or construction with classical correction compiled into a fully local noisy quantum circuit is also discussed by Selub, Bhardwaj and Lake~\cite[Section~1.2]{SBL2026}, citing Gottesman's textbook manuscript~\cite{Gottesman2024}. Shtanko and Sharma~\cite{SS2025} simulate local target circuits in every dimension without intermediate measurements, using a contracting nonunital channel with a unique nonmaximally mixed fixed point; the quantum refrigerator of Ben-Or, Gottesman and Hassidim~\cite{BGH2013} is a related reset-generating precedent. Our depolarizing noise is unital, and our elementary circuit has explicit resets.

Pattison, Haah and Preskill~\cite[Theorems~4.5 and~5.7]{PHP2026} prove universal fault-tolerant sampling on a planar device with a fixed depth-three periodic circuit away from a constant number of controlled qubits. Their processor receives an instruction stream during the run and uses a noiseless global clock. Here the target program is supplied entirely at time zero as part of a product preparation on a one-dimensional ring. Their manuscript was submitted on 18 November 2026; this is concurrent and independent work. Our rule also uses a fixed time-periodic schedule. The distinctions are one spatial dimension, a program supplied entirely in the initial state, and all classical control state, including the hierarchy's clocks, on noisy wires; like theirs, our layer schedule is applied noiselessly. Corollary~\ref{cor:fixed-channel} packages our periodic word into identical discrete updates by adding finitely many register banks. This is an elementary consequence of Theorem~\ref{thm:main}; its onsite bank shift acts on all $D_{\rm layer}Q_{\rm site}$ qubits at once and remains a noiseless timing structure. It does not remove the timing assumption in this comparison.


The classical foundation is G\'acs's initialized synchronous reliable automaton~\cite{Gacs1986}, rather than the stronger self-organizing asynchronous form~\cite{Gacs2001}. Tsirelson's local reliable memory~\cite{Cirelson1978}, G\'acs's two-dimensional construction~\cite{Gacs1989}, G\'acs--Reif's three-dimensional real-time array~\cite{GacsReif1988}, and Gray's account~\cite{Gray2001} place the hierarchy in its classical context. We use the explicit classical construction of Douglas and Mghirbi~\cite{DMNoise1D2026} and the one-dimensional local lemmas of their quantum memory~\cite{DMMemory2026}. The exact imported statements and all hypotheses used here are recorded in Sections~\ref{sec:inputs}--\ref{sec:code}.

Balasubramanian, Davydova and Lake~\cite[Section~1.3(3)]{BDL2026} formulate a quantum analogue of G\'acs's automaton as an open direction. Their two-dimensional toric-code memory has a reliable read-only instruction set and is not translation-invariant in spacetime. Drechsler, Oreg and Aharonov~\cite{DOA2026} relate CSS quantum memories to paired classical processes. For the 4D toric code they prove a Toom-rule memory threshold under local Pauli noise for a synchronous channel and for a periodically driven Lindbladian. They ask whether their framework helps in one dimension and leave a general time-independent lift from synchronous classical rules open. Our rule is discrete time with one fixed elementary period. Selub, Bhardwaj and Lake~\cite[Sections~1.2--1.3]{SBL2026} discuss local quantum computation with reliable classical processing and suggest G\'acs's self-simulation as a possible route to removing that assumption. Gidney and Bergamaschi~\cite[Theorem~1.2 and Section~11]{GB2025} obtain constant quantum space overhead with nonlocal noiseless classical control and quasi-polylogarithmic depth overhead in the range of their computation theorem, for compiled local target circuits. They ask for constant spatial overhead together with translation invariance. Our result supplies the fixed translation-invariant rule, with polynomial overhead. It does not establish their combined target.

The proof has three additional interfaces. Repeated classical bits are generally not quantum codewords, so a teleportation normalizer must first select the desired logical basis state. All parent registers, including program and control, must be encoded on the children's compiler tracks; the children's own processor runs on disjoint factors. Finally, a total classical decoder must agree with quantum recovery even on a bad intermediate block. Sections~\ref{sec:boundary}--\ref{sec:compiler} prove these channel statements. Section~\ref{sec:probability} obtains rank decay strong enough for a resource bound polynomial in logarithmic precision. Section~\ref{sec:finite} pays for the entire physical past cone on a finite ring, and Section~\ref{sec:proofs} proves Theorems A--C. Section~\ref{sec:general-noise} proves Theorem~\ref{thm:general}: the same circuit tolerates arbitrary sufficiently small local CPTP perturbations, and more generally local fault-path noise, with the same requested accuracy and polynomial resources. Its new ingredient is the dense-witness count for the one-dimensional hierarchy with overlapping long tests. The summation method is that of Aliferis, Gottesman and Preskill~\cite{AGP2006}, with related local-bath threshold methods of Terhal and Burkard~\cite{TB2005} and Aharonov, Kitaev and Preskill~\cite{AKP2006}; D\"unnweber, Styliaris and Trivedi~\cite[Definition~3]{DST2026} already use the same noise model in two dimensions. Section~\ref{sec:fixed-channel} proves the fixed-channel corollary and its reference-entangled storage guarantee through a horizon linear in ring size. Neither its ideal bank shift nor its finite-horizon guarantee supplies a time-independent continuous-time implementation or a guarantee after that horizon.
